\documentclass[10pt,a4paper]{amsart}
\usepackage[margin=19mm]{geometry}
\usepackage{amsmath,amssymb,mathtools}
\usepackage[scr=rsfs,cal=euler]{mathalfa}
\usepackage{xfrac}
\usepackage[unicode,hidelinks,bookmarks=false]{hyperref}
\newcommand{\ie}{i.e.\ }
\newcommand{\cf}{cf.\ }
\newcommand{\resp}{respectively\ }
\newcommand{\Subsection}{\subsection}
\providecommand{\nobreakdash}{\nobreak\hbox{-}}
\setlength{\emergencystretch}{2em}
\multlinegap0pt
\usepackage{bm}
\usepackage{bbm}
\usepackage{dsfont}
\usepackage{xspace}
\usepackage{cancel}
\usepackage{mathtools}
\usepackage{amsthm}
\makeatletter
\@ifundefined{theorem}{\newtheorem{theorem}{Theorem}}{}
\@ifundefined{lemma}{\newtheorem{lemma}[theorem]{Lemma}}{}
\makeatother
\newcommand{\CC}{\mathbb{C}}
\newcommand{\R}{\mathbb{R}}
\newcommand{\Z}{\mathbb{Z}}
\newcommand{\imag}{\operatorname{im}}
\newcommand{\real}{\operatorname{re}}
\title[Subconvexity Problem on $\operatorname{GL}\_3$ over number fi]{Subconvexity Problem on $\operatorname{GL}\_3$ over number fields: the twist aspect}
\author{Filippo Berta}
\date{Source: arXiv:2602.20095v1 (CC BY 4.0)}
\begin{document}
\maketitle

\noindent\emph{Source and licence.} Subconvexity Problem on $\operatorname{GL}_3$ over number fields: the twist aspect. Version arXiv:2602.20095v1 (CC BY 4.0), distributed under the Creative Commons Attribution 4.0 International licence, as recorded in the arXiv licence metadata for this exact version and shown on the abstract page https://arxiv.org/abs/2602.20095v1. Full rights notice: licenses/arxiv-2602.20095-CC-BY-4.0.txt.\par\smallskip

\noindent\textit{Editorial scope.} Counted unit: the Lemma whose source label is \texttt{eq:decaybesselforinert} in the article (source0.tex lines 776--783, source label \texttt{eq:decaybesselforinert}), delivered together with the article's own complete original proof of it (source0.tex lines 786--817, one \texttt{proof} environment of 32 lines). No mathematics is written, repaired or re-derived: statement and proof are the article's own text line for line. Required context retained uncounted: eq:inert (lines 767--767). Every definition, display and label of those lines is retained unchanged. Omitted from this package and not redistributed: 1--766, 768--775, 784--785, 818--1539. Nothing inside a retained excerpt is omitted or altered: every retained body is delivered exactly as the article's lines are written, so no omission receipt of any kind accompanies this package. Citations: the retained excerpts carry no citation command at all, so no bibliography entry of the article is transcribed and no reference is redistributed. Reference adaptation: none was needed and none was made; every reference inside the delivered unit resolves inside it, so no unresolved reference or citation is produced. Printed slips kept literal: no printed word, symbol, hypothesis or number of the counted statement or of its proof was altered. Layout and class: the article's own class is \texttt{amsart} with packages including graphicx, amsmath, amssymb, amsthm, enumerate, mathtools, tikz-cd, bbold, footmisc, mathrsfs, yfonts, geometry, hyperref; the wrapper is the lane's pinned \texttt{amsart} editorial wrapper at 10pt on A4, which restates the theorem-like environments the retained text needs and adds the article's own macro definitions that the retained text needs (\texttt{\textbackslash CC}, \texttt{\textbackslash R}, \texttt{\textbackslash Z}, \texttt{\textbackslash imag}, \texttt{\textbackslash real}), with extra wrapper packages (bm, bbm, dsfont, xspace, cancel, mathtools). The delivered numbering is the unit's own. Assets: no retained span contains a figure environment, a graphics-inclusion command, a PDF-page inclusion, a table or a TikZ picture; no image file is copied into this package, the package declares no asset dependency and no image file is redistributed with it. Licence: version arXiv:2602.20095v1 of this article is distributed under the Creative Commons Attribution 4.0 International licence, as recorded in the arXiv licence metadata for this exact version and shown on the abstract page https://arxiv.org/abs/2602.20095v1. A full rights notice is delivered at licenses/arxiv-2602.20095-CC-BY-4.0.txt. Characters: every retained excerpt is pure ASCII, so the delivered engine, XeLaTeX with its default Latin Modern text font, reports no missing character.
    We call a function $V\in C_c^{\infty}(\mathbb{C}^{\times})$ \emph{inert} if \[V(re^{i\theta}) \neq 0 \Rightarrow r \asymp 1\] and\begin{equation}\label{eq:inert}(r\partial_r)^n\partial_{\theta}^mV(re^{i\theta}) \ll_{n,m} 1.\end{equation}

\begin{lemma}
If $V \in C_c^{\infty}(\CC^{\times})$ is inert, then for any $m,n \in \Z_{\geq 0}$ and $A>0$ we have
\begin{equation}\label{eq:decaybesselforinert}(r\partial_r)^n\partial_{\theta}^mB_{\rho}[V](re^{i\theta}) \ll_{m,n,A}  \min(r^{2(1-Q)},r^{-A}).\end{equation}
If $V\in C_c^{\infty}(\R^{\times})$ is inert, then for any $m,n \in \Z_{\geq0 }$ and $A>0$ we have 
\begin{equation}\label{eq:decaybesselinert2}
    (x\partial_x)^nV(x) \ll_{n,A} \min(\vert x\vert^{1-Q},\vert x \vert^{-A}).
\end{equation}
\end{lemma}

\begin{proof}We prove only the complex case, the real one is easier. Since the real part $k$ also varies, we need to be careful in the use of the Stirling approximation formula: it says that for fixed $\sigma$, $\gamma$ will grow at most polynomially as  $\operatorname{im}(s)\to \infty$. The bound is in fact also polynomial in $k$ (once $\real(s)$ is fixed) as the following shows: By Stirling's approximation formula we have 
    \begin{align*}
    \frac{\Gamma_{\CC}\left(s+\overline{\alpha}+\frac{\vert{k}\vert}{2}\right)}{\Gamma_{\CC}\left(1-s+\alpha+\frac{\vert{k}\vert}{2}\right)} & \asymp e^{-2\real(s)}\pi^{1-2s}w_{s,\alpha,k}\left\vert{\frac{s+\alpha+\frac{\vert{k}\vert}{2}}{1-s+\alpha+\frac{\vert{k}\vert}{2}}}\right\vert^{\frac{\vert{k}\vert}{2}}\\&\times\left\vert{\left(1-s+\alpha+\frac{\vert{k}\vert}{2}\right)\left(s+\overline{\alpha}+\frac{\vert{k}\vert}{2}\right)}\right\vert^{\real(s)}\\&\times e^{-\operatorname{im}(s+\overline{\alpha})arg(s,\alpha,k)},
    \end{align*}
    where $w_{s,\alpha,k}=\left(\frac{1-s+\alpha+\frac{\vert{k}\vert}{2}}{s+\overline{\alpha}+\frac{\vert{k}\vert}{2}}\right)^{\frac{1}{2}-\real(\alpha)}$ and $arg(s,\alpha,k) = \operatorname{arg}(s+\overline{\alpha}+\frac{\vert{k}\vert}{2})+\operatorname{arg}(1-s+\alpha+\frac{\vert{k}\vert}{2})$.
  
    We need to make sure $arg(s,\alpha,k)$ is very small whenever $\imag(s)$ or $\vert{k}\vert$ is big. Assuming $\vert{k}\vert > 2(1-\real(s)+\real(\alpha))$, which eventually happens as the right-hand side can be fixed, we have \begin{align}\label{eq:arg}&arg(s,\alpha,k) =\nonumber\\& \arctan\left(\frac{\imag(s-\alpha)}{\real(s)+\real(\alpha)+\frac{\vert{k}\vert}{2}}\right)-\arctan\left(\frac{\imag(s-\alpha)}{1-\real(s)+\real(\alpha)+\frac{\vert{k}\vert}{2}}\right).\end{align}
    \begin{itemize}
        \item In the range $\vert{\imag(s-\alpha)}\vert> \vert{k}\vert$,  we can write \eqref{eq:arg} as 
        \begin{align*}&\sum_{n=0}^{\infty}\frac{(-1)^n}{(2n+1)}\left(\frac{(1-\real(s-\alpha)+\frac{\vert{k}\vert}{2})^{2n+1}-(\real(s+\alpha)+\frac{\vert{k}\vert}{2})^{2n+1}}{(\imag(s-\alpha))^{2n+1}}\right) \\ 
        &\ll \frac{1}{\vert{\imag(s-\alpha)}\vert}\sum_{n=0}^{\infty}\left\vert{\frac{k}{\imag(s-\alpha)}}\right\vert^{2n} = O(\frac{1}{\vert{\imag(s)}\vert}) \end{align*}
        \item The range $\vert{\imag(s-\alpha)}\vert<\frac{\vert{k}\vert}{2}$, we can write \eqref{eq:arg} as 
        \begin{align*}
            &\sum_{n=0}^{\infty}\frac{(-1)^n}{(2n+1)}\left(\frac{(\imag(s-\alpha))^{2n+1}}{(1-\real(s-\alpha)+\frac{\vert{k}\vert}{2})^{2n+1}}-\frac{(\imag(s-\alpha))^{2n+1}}{(\real(s+\alpha)+\frac{\vert{k}\vert}{2})^{2n+1}}\right)\\
            &\ll \frac{1}{\vert{k}\vert}\sum_{n=0}^{\infty}\left\vert{\frac{\imag(s-\alpha)}{k/2}}\right\vert^{2n+1} = O(\frac{1}{\vert{k}\vert}) 
        \end{align*}
        \item Finally in the range $\frac{\vert{k}\vert}{2}\leq \vert{\imag(s-\alpha)}\vert\leq \vert{k}\vert$ we use the fact that $\arctan$ is Lipschitz continuous to say that \eqref{eq:arg} is less or equal to
        \begin{align*}
            &\ll\left\vert{\frac{\imag(s-\alpha)}{\real(s)+\real(\alpha)+\frac{\vert{k}\vert}{2}}-\frac{\imag(s-\alpha)}{1-\real(s)+\real(\alpha)+\frac{\vert{k}\vert}{2}}}\right\vert
            \\& \ll \left\vert{\frac{\imag(s-\alpha)}{k^2}}\right\vert \ll \frac{1}{\vert{\imag(s-\alpha)}\vert}
        \end{align*}
         \end{itemize}
        In any case we see that $e^{-\imag(s-\alpha)arg(s,\alpha,k)} = O(1)$. We deduce that for fixed $\real(s)$ there exists $C>0$ (possibly depending on $\real(s)$ and the set of archimedean parameters of $\pi$) so that 
        \[\gamma(1-s,\rho_k,\psi) \ll \left((1+\vert{k}\vert)(1+\vert{t}\vert)\right)^C. \]
        Shifting the countour of integration to $Q$ we have 
        \begin{align*}
            &B_{\rho}[V](re^{i\theta}) \\ &\ll r^{2-2Q}\sum_{k\in \Z}\int_{-\infty}^{\infty}\left\vert{\gamma(1-Q-it,\rho_k,\psi) \mathcal{M}(V,k)(-2(Q+it))}\right\vert \ \mathrm{d}t \\ &\ll r^{2(1-Q)}\sum_{k\in\Z}\int_{-\infty}^{\infty}\left((1+\vert{k}\vert)(1+\vert{t}\vert)\right)^C\vert{\mathcal{M}(V,k)(-2(Q+it))}\vert \ \mathrm{d}t \ll r^{2(1-Q)},
        \end{align*}
        where the last inequality follows from \eqref{eq:inert}.
        Similarly, for the bound $B_{\pi}[V](re^{i\theta}) \ll_A r^{-A}$ one shifts the countour to $\real(s) = \frac{A}{2}$.  
       
    \end{proof}

\end{document}
